% АВТОГЕНЕРАЦИЯ: build_references.py <- references.bib, порядок первого цитирования в article.tex
\ntvref{kulkarni2015}{Kulkarni A., Callan J. Selective search: Efficient and effective search of large textual collections // ACM Transactions on Information Systems. 2015. V. 33. N 4. P. 17:1–17:33. \url{https://doi.org/10.1145/2738035}}{Kulkarni A., Callan J. Selective search: Efficient and effective search of large textual collections. \textit{ACM Transactions on Information Systems}, 2015, vol.~33, no.~4, pp.~17:1–17:33. \url{https://doi.org/10.1145/2738035}}
\ntvref{dai2016}{Dai Z., Xiong C., Callan J. Query-biased partitioning for selective search // Proceedings of the 25th ACM International Conference on Information and Knowledge Management (CIKM). 2016. P. 1119–1128. \url{https://doi.org/10.1145/2983323.2983706}}{Dai Z., Xiong C., Callan J. Query-biased partitioning for selective search. \textit{Proceedings of the 25th ACM International Conference on Information and Knowledge Management (CIKM)}, 2016, pp.~1119–1128. \url{https://doi.org/10.1145/2983323.2983706}}
\ntvref{risvik2013}{Risvik K.M., Chilimbi T., Tan H. et al. Maguro, a system for indexing and searching over very large text collections // Proceedings of the 6th ACM International Conference on Web Search and Data Mining (WSDM). 2013. P. 727–736. \url{https://doi.org/10.1145/2433396.2433486}}{Risvik K.M., Chilimbi T., Tan H. et al. Maguro, a system for indexing and searching over very large text collections. \textit{Proceedings of the 6th ACM International Conference on Web Search and Data Mining (WSDM)}, 2013, pp.~727–736. \url{https://doi.org/10.1145/2433396.2433486}}
\ntvref{moffat2007}{Moffat A., Webber W., Zobel J., Baeza-Yates R. A pipelined architecture for distributed text query evaluation // Information Retrieval. 2007. V. 10. N 3. P. 205–231. \url{https://doi.org/10.1007/s10791-006-9014-4}}{Moffat A., Webber W., Zobel J., Baeza-Yates R. A pipelined architecture for distributed text query evaluation. \textit{Information Retrieval}, 2007, vol.~10, no.~3, pp.~205–231. \url{https://doi.org/10.1007/s10791-006-9014-4}}
\ntvref{moffat2006}{Moffat A., Webber W., Zobel J. Load balancing for term-distributed parallel retrieval // Proceedings of the 29th Annual International ACM SIGIR Conference on Research and Development in Information Retrieval. 2006. P. 348–355. \url{https://doi.org/10.1145/1148170.1148232}}{Moffat A., Webber W., Zobel J. Load balancing for term-distributed parallel retrieval. \textit{Proceedings of the 29th Annual International ACM SIGIR Conference on Research and Development in Information Retrieval}, 2006, pp.~348–355. \url{https://doi.org/10.1145/1148170.1148232}}
\ntvref{bouma2009}{Bouma G. Normalized (pointwise) mutual information in collocation extraction // Proceedings of the Biennial GSCL Conference. 2009. P. 31–40.}{Bouma G. Normalized (pointwise) mutual information in collocation extraction. \textit{Proceedings of the Biennial GSCL Conference}, 2009, pp.~31–40.}
\ntvref{williams2022}{Williams C.K.I. On suspicious coincidences and pointwise mutual information // Neural Computation. 2022. V. 34. N 10. P. 2037–2046. \url{https://doi.org/10.1162/neco_a_01533}}{Williams C.K.I. On suspicious coincidences and pointwise mutual information. \textit{Neural Computation}, 2022, vol.~34, no.~10, pp.~2037–2046. \url{https://doi.org/10.1162/neco_a_01533}}
\ntvref{dunning1993}{Dunning T. Accurate methods for the statistics of surprise and coincidence // Computational Linguistics. 1993. V. 19. N 1. P. 61–74. \url{https://aclanthology.org/J93-1003}}{Dunning T. Accurate methods for the statistics of surprise and coincidence. \textit{Computational Linguistics}, 1993, vol.~19, no.~1, pp.~61–74. \url{https://aclanthology.org/J93-1003}}
\ntvref{robertson2009}{Robertson S., Zaragoza H. The probabilistic relevance framework: {BM25} and beyond // Foundations and Trends in Information Retrieval. 2009. V. 3. N 4. P. 333–389. \url{https://doi.org/10.1561/1500000019}}{Robertson S., Zaragoza H. The probabilistic relevance framework: {BM25} and beyond. \textit{Foundations and Trends in Information Retrieval}, 2009, vol.~3, no.~4, pp.~333–389. \url{https://doi.org/10.1561/1500000019}}
\ntvref{newman2004}{Newman M.E.J., Girvan M. Finding and evaluating community structure in networks // Physical Review E. 2004. V. 69. N 2. P. 026113. \url{https://doi.org/10.1103/PhysRevE.69.026113}}{Newman M.E.J., Girvan M. Finding and evaluating community structure in networks. \textit{Physical Review E}, 2004, vol.~69, no.~2, pp.~026113. \url{https://doi.org/10.1103/PhysRevE.69.026113}}
\ntvref{reichardt2006}{Reichardt J., Bornholdt S. Statistical mechanics of community detection // Physical Review E. 2006. V. 74. N 1. P. 016110. \url{https://doi.org/10.1103/PhysRevE.74.016110}}{Reichardt J., Bornholdt S. Statistical mechanics of community detection. \textit{Physical Review E}, 2006, vol.~74, no.~1, pp.~016110. \url{https://doi.org/10.1103/PhysRevE.74.016110}}
\ntvref{traag2019}{Traag V.A., Waltman L., van Eck N.J. From {Louvain} to {Leiden}: guaranteeing well-connected communities // Scientific Reports. 2019. V. 9. N 1. P. 5233. \url{https://doi.org/10.1038/s41598-019-41695-z}}{Traag V.A., Waltman L., van Eck N.J. From {Louvain} to {Leiden}: guaranteeing well-connected communities. \textit{Scientific Reports}, 2019, vol.~9, no.~1, pp.~5233. \url{https://doi.org/10.1038/s41598-019-41695-z}}
\ntvref{traag2011}{Traag V.A., Van Dooren P., Nesterov Y. Narrow scope for resolution-limit-free community detection // Physical Review E. 2011. V. 84. N 1. P. 016114. \url{https://doi.org/10.1103/PhysRevE.84.016114}}{Traag V.A., Van Dooren P., Nesterov Y. Narrow scope for resolution-limit-free community detection. \textit{Physical Review E}, 2011, vol.~84, no.~1, pp.~016114. \url{https://doi.org/10.1103/PhysRevE.84.016114}}
\ntvref{rosvall2008}{Rosvall M., Bergstrom C.T. Maps of random walks on complex networks reveal community structure // Proceedings of the National Academy of Sciences. 2008. V. 105. N 4. P. 1118–1123. \url{https://doi.org/10.1073/pnas.0706851105}}{Rosvall M., Bergstrom C.T. Maps of random walks on complex networks reveal community structure. \textit{Proceedings of the National Academy of Sciences}, 2008, vol.~105, no.~4, pp.~1118–1123. \url{https://doi.org/10.1073/pnas.0706851105}}
\ntvref{karypis1998}{Karypis G., Kumar V. A fast and high quality multilevel scheme for partitioning irregular graphs // SIAM Journal on Scientific Computing. 1998. V. 20. N 1. P. 359–392. \url{https://doi.org/10.1137/S1064827595287997}}{Karypis G., Kumar V. A fast and high quality multilevel scheme for partitioning irregular graphs. \textit{SIAM Journal on Scientific Computing}, 1998, vol.~20, no.~1, pp.~359–392. \url{https://doi.org/10.1137/S1064827595287997}}
\ntvref{chvatal1979}{Chv{\'a}tal V. A greedy heuristic for the set-covering problem // Mathematics of Operations Research. 1979. V. 4. N 3. P. 233–235. \url{https://doi.org/10.1287/moor.4.3.233}}{Chv{\'a}tal V. A greedy heuristic for the set-covering problem. \textit{Mathematics of Operations Research}, 1979, vol.~4, no.~3, pp.~233–235. \url{https://doi.org/10.1287/moor.4.3.233}}
\ntvref{bajaj2016}{Bajaj P., Campos D., Craswell N. et al. {MS MARCO}: A human generated {MAchine Reading COmprehension} dataset // arXiv preprint arXiv:1611.09268. 2016. \url{https://doi.org/10.48550/arXiv.1611.09268}}{Bajaj P., Campos D., Craswell N. et al. {MS MARCO}: A human generated {MAchine Reading COmprehension} dataset. \textit{arXiv preprint arXiv:1611.09268}, 2016. \url{https://doi.org/10.48550/arXiv.1611.09268}}
\ntvref{kwiatkowski2019}{Kwiatkowski T., Palomaki J., Redfield O. et al. Natural Questions: a benchmark for question answering research // Transactions of the Association for Computational Linguistics. 2019. V. 7. P. 453–466. \url{https://doi.org/10.1162/tacl_a_00276}}{Kwiatkowski T., Palomaki J., Redfield O. et al. Natural Questions: a benchmark for question answering research. \textit{Transactions of the Association for Computational Linguistics}, 2019, vol.~7, pp.~453–466. \url{https://doi.org/10.1162/tacl_a_00276}}
